Your Agent is Underperforming:

The Case for Agent Performance Improvement Plans

Anonymous Submission --- Double-Blind Review

Author information withheld per review guidelines

Abstract

Production AI agents frequently degrade after deployment, yet structured remediation culture is largely absent. Empirical studies confirm most deployments fail or underdeliver (Pan et al., 2025), and simulation-based research projects that unchecked behavioral drift reduces task success by 42\% (Rath, 2026). We introduce the Agent Performance Improvement Plan (APIP): a lifecycle protocol adapted from organizational management that governs the remediation of underperforming agents. Drawing on human PIP structure and Kirkpatrick's four-level evaluation framework, the APIP defines trigger thresholds, a four-category failure mode taxonomy, a tiered intervention hierarchy, a bounded remediation window with check-ins, and exit criteria. A discrete-event simulation across all four failure modes shows that tier-mismatched ad-hoc remediation extends time-to-resolution by 2.1$\times$, validating the APIP's cause attribution requirement. We formalize a declarative schema meeting AI governance traceability standards (Kolt et al., FAccT 2024; EU AI Act), and ground extension to multi-agent meshes---where a locally optimal new agent degrades peer performance through context disruption (Kim et al., 2025)---in Event System Theory and team coordination research.

